% 19 November 2026. Drafted for Edward Baker with OpenAI Codex assistance.
% Helper explicitly dispatched as gpt-6.1-sol, reasoning effort ultra;
% independent serving-backend attestation was not available to the helper.
\section{The forced residual and the joint gap readout}
\label{sec:residual-transport}

We next distinguish an exact heat evolution from its explicit arithmetic
approximation.  This distinction determines which error may be
propagated.  All calculations in this section use a fixed cutoff, or
one cell on which the cutoff is constant.  They are elementary
consequences of the source representation and the physical equation
$H_t=-H_{zz}$.  They supply an interface for a quantitative estimate;
no such new canonical estimate is asserted here.

\subsection{The exact term evolves by heat; the explicit term has a defect}

For the exact Gaussian term \eqref{as:exact-term}, differentiation and
integration by parts give
\begin{align}
 \partial_t r_n(t,s)
  &=\frac1{2\sqrt{\pi t}}\int_{\R}
       v\,r_{0,n}'(s+\sqrt t\,v)e^{-v^2}\,\dd v\notag\\
  &=\frac1{4\sqrt\pi}\int_{\R}
       r_{0,n}''(s+\sqrt t\,v)e^{-v^2}\,\dd v
    =\frac14\partial_s^2r_n(t,s).
 \label{rt:exact-term-heat}
\end{align}
The sub-Gaussian growth described after \eqref{as:exact-term} kills
the boundary term and justifies the differentiations on compact
nonreal $s$-sets and positive-time intervals.  The same calculation
applies to the exact remainder where its Gaussian formula is used.

There is no corresponding exact heat assertion for $p_n$.  Suppress
the index $n$, and write
\begin{equation}
 a=\alpha-\log n,\qquad b=\alpha',\qquad c=\alpha'',\qquad
 p=\exp(\ell_n+ta^2/4).
 \label{rt:abc}
\end{equation}
These are functions of $s$; explicitly
\[
 b(s)=-\frac1{2s^2}-\frac1{(s-1)^2}+\frac1{2s},\qquad
 c(s)=\frac1{s^3}+\frac2{(s-1)^3}-\frac1{2s^2}.
\]
Differentiating this actual exponential, rather than a remainder
symbol, gives
\begin{align}
 \frac{p_t}{p}&=\frac{a^2}{4},\qquad
 \frac{p_s}{p}=a+\frac t2ab,\notag\\
 \frac{p_{ss}}p&=\left(a+\frac t2ab\right)^2
                   +b+\frac t2(b^2+ac).
 \label{rt:pn-jets}
\end{align}
Subtracting one fourth of the last expression from the first gives
the exact defect
\begin{equation}
 \mathcal R_n(t,s):=\left(\partial_t-\frac14\partial_s^2\right)p_n(t,s)
 =-\frac{p_n(t,s)}4
   \left\{b+ta^2b+\frac t2(b^2+ac)
                       +\frac{t^2}{4}a^2b^2\right\}.
 \label{rt:term-defect}
\end{equation}
For the additive term error $d_n=r_n-p_n$ this implies
\begin{equation}
 \left(\partial_t-\frac14\partial_s^2\right)d_n=-\mathcal R_n.
 \label{rt:term-error-equation}
\end{equation}
The term $b$ in \eqref{rt:term-defect} is present even at $t=0$:
the explicit expression has replaced both the gamma factor and its
Gaussian evolution by an approximation.  A small value of the defect
is not a bound for the solution of \eqref{rt:term-error-equation}
until its data and domain are specified.

Because $\partial_z^2p_n(t,s_\pm(z))=-p_{n,ss}(t,s_\pm(z))/4$,
the physical fixed-$N$ reference has defect
\begin{equation}
 \mathcal R_N(t,z):=\partial_tP_N(t,z)+\partial_z^2P_N(t,z)
   =\sum_{n=1}^N\{\mathcal R_n(t,s_+(z))
                      +\mathcal R_n(t,s_-(z))\}.
 \label{rt:physical-defect}
\end{equation}
Define the actual error by $E=H-P_N$.  Then, exactly,
\begin{equation}
 E_t+E_{zz}=-\mathcal R_N.
 \label{rt:physical-error}
\end{equation}
For backward physical time $u=t_*-t$, with every quantity evaluated
at $t_*-u$, this becomes the forward forced heat equation
\begin{equation}
 E_u=E_{zz}+\mathcal R_N.
 \label{rt:backward-error}
\end{equation}
The sign and coefficient here follow from both changes of variables:
$s_\pm=(1\pm\ii z)/2$ and $u=t_*-t$.  There is no missing factor of
four in the physical equation.

Using the sum of the exact $r_n$ terms as the reference would instead
give zero defect, by \eqref{rt:exact-term-heat}; its full remainder
would still need to be controlled.  Using the explicit $p_n$ terms
gives a directly differentiable formula and the known forcing
\eqref{rt:physical-defect}.  These are two different acquisition
interfaces, not a reason to discard the remainder in either one.

\subsection{What is required to propagate the residual}

Equation \eqref{rt:backward-error} on a bounded spatial neighborhood
does not determine an effective error enclosure there from a local
initial bound alone.  For example, for $A>0$ the entire functions
\[
 E_K(u,z)=\varepsilon\exp\{K(z-A)+K^2u\},\qquad K>0,
\]
satisfy $E_u=E_{zz}$ and
$\sup_{-A\le x\le A}|E_K(0,x)|\le\varepsilon$, whereas
$E_K(u,A)=\varepsilon e^{K^2u}$.  Thus the same local initial upper
bound permits arbitrarily large later values.  The uncontrolled
growth outside the interval is exactly what prevents a uniform
conclusion.  This example concerns the information in a local heat
bound, not the canonical source.

One possible propagation statement is the following deliberately
conditional version.  Its strong global hypotheses are displayed
because they have not been acquired for the arithmetic reference.

\begin{proposition}[A whole-strip residual bound]
\label{rt:duhamel}
Let $0\le u\le U$ and $h>0$.  Suppose that $E$ and $R$ are
holomorphic in $z$ on a neighborhood of the closed horizontal strip
$|\Im z|\le h$, continuous in time, and that $E$ is a classical
solution of $E_u=E_{zz}+R$ there.  Suppose that $E$ is uniformly
bounded on this whole time-strip, that
\begin{equation}
 \sup_{x\in\R,\,|\eta|\le h}|E(0,x+\ii\eta)|\le\varepsilon_0,
 \qquad
 \sup_{x\in\R,\,|\eta|\le h}|R(v,x+\ii\eta)|\le r(v),
 \label{rt:global-data}
\end{equation}
where $r$ is integrable.  Assume the usual continuity needed to take
these bounded classical solutions to their initial data.  Then
\begin{equation}
 \sup_{x\in\R,\,|\eta|\le h}|E(u,x+\ii\eta)|
       \le\varepsilon(u):=\varepsilon_0+\int_0^u r(v)\,\dd v.
 \label{rt:duhamel-bound}
\end{equation}
For $0<h'<h$, Cauchy's estimate also gives
\begin{equation}
 \sup_{x\in\R,\,|\eta|\le h'}|E_{zz}(u,x+\ii\eta)|
       \le\frac{2\varepsilon(u)}{(h-h')^2}.
 \label{rt:duhamel-jets}
\end{equation}
\end{proposition}

\begin{proof}
For each fixed $\eta$, put
$G_u(q)=(4\pi u)^{-1/2}\exp(-q^2/(4u))$.  The bounded whole-line
heat solution has the Duhamel representation
\begin{align*}
 E(u,x+\ii\eta)={}&\int_{\R}G_u(x-\xi)E(0,\xi+\ii\eta)\,\dd\xi\\
 &+\int_0^u\int_{\R}G_{u-v}(x-\xi)
                      R(v,\xi+\ii\eta)\,\dd\xi\,\dd v.
\end{align*}
The heat kernel has integral one and is nonnegative on the real line.
Taking absolute values proves \eqref{rt:duhamel-bound}, uniformly in
$\eta$.  The representation follows by the fundamental solution and
uniqueness in the bounded classical class; it can also be verified by
regularizing in time and passing to the bounded initial data.  Apply
Cauchy's formula on disks of radius smaller than $h-h'$ and then take
their radius to $h-h'$ to obtain \eqref{rt:duhamel-jets}.
\end{proof}

The hypotheses include every real $x$, not just the part of the strip
where a zero pair is enclosed.  In particular this proposition is not
automatically applicable to the branch-defined reference $P_N$ of
\eqref{as:PN}, which was defined on a right half-plane.  A genuine use
must provide an admissible extension and the whole-line bounds, or use
a bounded-domain heat problem with certified lateral boundary data.
If a Gaussian integral is truncated, its tails must be paid as well.
For example, a global bound $M$ on its initial integrand pays at most
$M\operatorname{erfc}(A/(2\sqrt u))$ outside $|x-\xi|\le A$;
without such a bound, the displayed tail allowance has no premise.

Normalization does not remove this obligation.  If $H=\mathcal B F$
with a nonvanishing analytic factor, then $F$ satisfies
\begin{equation}
 F_t+F_{zz}+2\frac{\mathcal B_z}{\mathcal B}F_z
       +\frac{\mathcal B_t+\mathcal B_{zz}}{\mathcal B}F=0.
 \label{rt:normalized-equation}
\end{equation}
Thus a bound for the normalized error $H/\mathcal B-f$ cannot be
propagated by the plain unit-coefficient heat kernel while ignoring
the two additional coefficients.  Moreover a small absolute physical
error can be large relative to $H$ near a zero.  The contour denominator
needed next is a separate quantitative datum.

\subsection{A two-zero contour retains a joint error}

Let $\Gamma$ be a positively oriented piecewise smooth simple closed
contour, fixed over the time cell.  Assume that $H$ and the reference
$P$ are holomorphic on a neighborhood of its closed interior, nonzero
on a neighborhood of $\Gamma$, and each has exactly two zeros inside,
counted with multiplicity.  In this subsection
$P$ can be $P_N$ or another explicitly specified analytic reference.
The exact heat law is imposed on $H$, while
\begin{equation}
 R=P_t+P_{zz},\qquad E=H-P
 \label{rt:general-reference}
\end{equation}
is retained for $P$.  Counts and contour clearance must hold throughout
the cell; they are not inferred from a sampled value.

For $F=H$ or $P$, define
\begin{equation}
 A_j^F=\frac1{2\pi\ii}\oint_\Gamma z^j\frac{F_z}{F}\,\dd z,
 \qquad c_F=\frac{A_1^F}{2},\qquad
 D_F=\frac{A_2^F}{2}-\frac{(A_1^F)^2}{4}.
 \label{rt:pair-moments}
\end{equation}
If the two zeros are $a,b$, the argument principle gives
$c_F=(a+b)/2$ and $D_F=(b-a)^2/4$.  These contour expressions remain
defined when the two occurrences merge.  For a conjugation-invariant
pair, $D_F$ is real; it is positive for distinct real zeros and
negative for a distinct nonreal conjugate pair.

Differentiate the logarithmic derivative of $H$ and integrate by parts:
\begin{equation}
 (A_j^H)_t=\frac{j}{2\pi\ii}\oint_\Gamma
                          z^{j-1}\frac{H_{zz}}H\,\dd z,
 \qquad
 (D_H)_t=\frac1{2\pi\ii}\oint_\Gamma
                          (z-c_H)\frac{H_{zz}}H\,\dd z.
 \label{rt:exact-gap-readout}
\end{equation}
For example,
$\partial_t(H_z/H)=\partial_z(H_t/H)
=-\partial_z(H_{zz}/H)$, which proves the first formula.  Substitution
in the definition of $D_H$ proves the second.  This is an exact readout
of the pair law, including at a merger, rather than a new exclusion
criterion.

On the cleared contour, elementary algebra gives
\begin{equation}
 \frac{H_{zz}}H-\frac{P_{zz}}P
       =\frac{P E_{zz}-P_{zz}E}{P(P+E)}.
 \label{rt:joint-quotient}
\end{equation}
The two error terms in the numerator share the same residual.  For a
spatially constant multiplier $E=\lambda(t)P$, the numerator vanishes
exactly.  Separate bounds for $P E_{zz}$ and $P_{zz}E$ lose that
dependence.  This identity identifies a possible place to retain
arithmetic cancellation, but it proves no bound for that cancellation
in the canonical source.

There are two further payments.  The actual center in
\eqref{rt:exact-gap-readout} is $c_H$, not $c_P$.  Also $P_{zz}/P$
alone does not give the time motion of the zeros of $P$.  Repeating
the preceding calculation with $P_t=-P_{zz}+R$ gives
\begin{equation}
 (D_P)_t=\frac1{2\pi\ii}\oint_\Gamma
                   (z-c_P)\left(\frac{P_{zz}}P-\frac RP\right)\,\dd z.
 \label{rt:reference-gap-readout}
\end{equation}
Subtracting it from \eqref{rt:exact-gap-readout} yields the full identity
\begin{align}
 (D_H-D_P)_t=\frac1{2\pi\ii}\oint_\Gamma
 \bigg\{&(z-c_H)\frac{P E_{zz}-P_{zz}E}{P(P+E)}
       +(c_P-c_H)\frac{P_{zz}}P\notag\\
       &+(z-c_P)\frac RP\bigg\}\,\dd z.
 \label{rt:full-gap-error}
\end{align}
In particular, paying only \eqref{rt:joint-quotient} and calling its
integral the error in $(D_P)_t$ would omit both a center error and the
known heat defect.

The center can itself be enclosed without differentiating a value
error.  Suppose $|E/P|\le q<1$ on a neighborhood of $\Gamma$.
The convergent branch $\log(1+E/P)$ is then single-valued there.
Integration by parts along the contour gives
\begin{equation}
 c_H-c_P=-\frac1{4\pi\ii}\oint_\Gamma
                          \log(1+E/P)\,\dd z,
 \qquad
 |c_H-c_P|\le\frac{|\Gamma|}{4\pi}\{-\log(1-q)\}.
 \label{rt:center-bound}
\end{equation}
Here $|\Gamma|$ is the contour length.  No logarithm is asserted through
the interior zeros of $H$ or $P$; only the collar of the contour is used.
The same smallness condition, together with the two-zero count for $P$,
also supplies the count for $H$ by Rouch\'e's theorem.

For a concrete bound on \eqref{rt:full-gap-error}, set on each slice
\begin{gather*}
 m=\inf_\Gamma|P|,\quad \varepsilon=\sup_\Gamma|E|<m,\quad
 J=\sup_\Gamma|P E_{zz}-P_{zz}E|,\\
 Q=\sup_\Gamma|P_{zz}/P|,\quad W=\sup_\Gamma|R/P|,\quad
 B_H=\sup_\Gamma|z-c_H|,\quad B_P=\sup_\Gamma|z-c_P|.
\end{gather*}
Then
\begin{equation}
 |(D_H-D_P)_t|\le\mathcal A(t):=
 \frac{|\Gamma|}{2\pi}
 \left\{\frac{B_HJ}{m(m-\varepsilon)}
              +|c_H-c_P|Q+B_PW\right\}.
 \label{rt:gap-error-bound}
\end{equation}
The $J$ bound may be obtained directly from a coupled enclosure, or
estimated less sharply from paid separate jets.  The latter is a valid
fallback, with no promised gain from correlation.  Every denominator
floor is displayed.  Integrating in backward time gives the conditional
comparison
\begin{equation}
 |D_H(u)-D_P(u)|\le |D_H(0)-D_P(0)|
                      +\int_0^u\mathcal A(t_*-v)\,\dd v.
 \label{rt:integrated-gap-error}
\end{equation}
To use this for continuation, the right side must be smaller than an
available positive lower bound for $D_P(u)$, with the counts, reality
symmetry, and all-time contour clearance retained.  Alternatively one
can bound the exact drift in \eqref{rt:exact-gap-readout} directly and
feed its integrated shrinking allowance to Section~\ref{sec:continuation}.
Neither use is supplied by a small residual norm in isolation.

\subsection{The center of a corrected approximation must also change}

The leading Riemann--Siegel correction gives two different value
approximations.  Schematically, with the source normalizer,
\begin{equation}
 \frac{H_t}{\mathcal B_t}=f_{t,N}+e_0,
 \qquad
 \frac{H_t}{\mathcal B_t}
       =\left(f_{t,N}-\frac{C_t}{\mathcal B_t}\right)+e_1.
 \label{rt:corrected-centers}
\end{equation}
The published bound for $e_1$ replaces the omitted-correction budget
$e_{C,0}$ by the smaller retained-correction budget $e_C$
\cite[\S6, (69)--(70)]{Polymath2019}.  A lower bound $|f_{t,N}|\ge L$
therefore does not automatically consume the $e_1$ allowance.  It needs
a lower bound for $|f_{t,N}-C_t/\mathcal B_t|$.  The unconditional
reverse-triangle transfer pays
\begin{equation}
 \left|f_{t,N}-\frac{C_t}{\mathcal B_t}\right|
          \ge L-\left|\frac{C_t}{\mathcal B_t}\right|,
 \label{rt:correction-payment}
\end{equation}
which can spend the margin recovered in the remainder.  This was the
applicability issue for the triangle estimate examined in
\cite{NewmanCertificateAudit}.  It is scoped to that estimate and its
information: it does not exclude phase-sensitive lower bounds or a
different arithmetic reorganization.

There is an additional analytic issue for a contour consumer.  A
displayed correction written in terms of $x,y$, a real part, and a
floor-valued cutoff is a value formula.  Before taking its complex
derivatives, one must establish a holomorphic representative on the
chosen neighborhood or use a different analytic reference with a
separately paid correction.  Simply assigning a smaller value error
to $P_N-C_t$ does not establish that reference's analyticity or heat
defect.  If such a reference is supplied, all of
\eqref{rt:physical-defect}--\eqref{rt:gap-error-bound} must use that
same reference, its center, and its actual defect.

Arithmetic mollifiers illustrate another legitimate cancellation that
requires a transfer.  For a finite Dirichlet polynomial
$A(s)=\sum_{n\le N}a_n n^{-s}$ and
$m(s)=\sum_{d\le D}\lambda_d d^{-s}$, exact multiplication gives
\begin{equation}
 m(s)A(s)=\sum_{r\le DN}
           \left(\sum_{dn=r}\lambda_d a_n\right)r^{-s}.
 \label{rt:mollifier-convolution}
\end{equation}
One can combine the coefficients in the inner sum before taking
absolute values.  If a full calculation establishes $|m f|\ge Q>0$
and $|m|\le M$, then $m\ne0$ and
\begin{equation}
 |f|\ge Q/M.
 \label{rt:mollifier-transfer}
\end{equation}
Gomila's finite-sum binding applies this type of arithmetic
reorganization to the two pieces of the normalized approximation and
returns a bound for the original $f_t$
\cite{Gomila2026Binding,VCDCMDiscussion12}.  The convolution and division
are mathematical steps with their own positive-margin hypotheses.

For a moving gap, multiplication is not a free change of the physical
equation.  If $P$ has defect $R$ and $m=m(t,z)$ is a nonvanishing
analytic factor on the relevant region, then
\begin{equation}
 (mP)_t+(mP)_{zz}
      =mR+(m_t+m_{zz})P+2m_zP_z.
 \label{rt:factor-defect}
\end{equation}
The factor preserves the enclosed zero divisor when it has no zeros
there, but it changes the heat defect and the second-logarithmic
quotient.  The zero dynamics are invariant only when these extra terms
are retained in the full readout.  Thus the static success of a
mollifier lower bound does not imply that the multiplied field solves
the same Newman heat equation.

Finally the generic residual interface must preserve actual collisions.
For the exact control $H=P=z^2-2(t-T)$ one has $E=R=0$, while the two
real zeros merge at $t=T$.  Here $D=2(t-T)$, so an integrated backward
shrinking allowance reaches the initial $D$ exactly at the merger;
the strict continuation margin fails.  Perfect approximation has not
protected a colliding reference.  The more detailed premise screen in
Section~\ref{sec:program} distinguishes the complete canonical
theta source from the colliding source controls.  No generic error
identity above supplies that distinction.  The open arithmetic task is
to acquire a complete, source-specific bound at the final gap or strip
consumer, and to compare it with direct evaluation using the same
information and error payments.
